\documentclass[11pt,letterpaper]{article}

\usepackage[margin=1in]{geometry}
\usepackage{graphicx}
\usepackage{booktabs}
\usepackage{tabularx}
\usepackage{listings}
\usepackage[hidelinks]{hyperref}

\lstdefinestyle{evidence}{
    basicstyle=\ttfamily\scriptsize,
    frame=single,
    breaklines=true,
    breakatwhitespace=false,
    columns=fullflexible,
    keepspaces=true,
    showstringspaces=false,
    captionpos=b
}
% Support builds from either the repository root or documentation/.
\IfFileExists{evidence/run-59ewk1/01-silver.txt}{
    \newcommand{\evidenceroot}{evidence/run-59ewk1}
}{
    \newcommand{\evidenceroot}{documentation/evidence/run-59ewk1}
}

% Based on the Report requirements in Team_Project (4).pdf.
% Main report: at most eight pages. An appendix is permitted if needed.
% Submit the report PDF with the final code as a zip file.
% Comments indicate the sections still to be completed.

\title{CPS 853 Team Project Report}
\author{} % Add team members' names.
\date{}

\begin{document}

\maketitle

\section{Design}
% 2 marks.

\subsection{Architecture Diagram}
% Insert the architecture diagram and explain the main components/data flow.
% Example figure structure (uncomment after adding the image):
% \begin{figure}[htbp]
%     \centering
%     \includegraphics[width=\linewidth]{architecture.pdf}
%     \caption{System architecture.}
%     \label{fig:architecture}
% \end{figure}

\subsection{Database Choice and Alternative}
% Compare the selected database with one alternative and justify the choice.

\subsection{Changes Since the Proposal}
% Describe any changes and give the reason for each.

\section{Evidence That All Four Parts Work}
% 3 marks. Include sample outputs, stage counts, quality-check results,
% and the three answers with measured database response times.

\subsection{Part 1: Store and Batch}
The Spark batch job reads the raw hourly PM2.5 CSV, checks required fields,
numeric values, timestamps, the 2024 date range, and duplicate measurement keys,
then writes accepted observations and rejected records to separate Parquet
datasets. In this captured run, all 8,139,551 input rows passed validation.
The subsequent exclusion of state codes 72 and 80 is a separate filtering step;
the validation count is therefore not the final Silver row count.

\noindent\begin{minipage}{\linewidth}
\lstinputlisting[
    style=evidence,firstline=2,lastline=13,
    caption={Captured Silver processing: CSV validation and Parquet outputs.},
    label={lst:batch-silver}
]{\evidenceroot/01-silver.txt}
\end{minipage}


State and county dimensions are loaded from Silver before publishing batch
aggregates. The captured outputs show 51 state entries and 587 counties,
followed by 51 state-year aggregates and 205,814 county-day aggregates saved
to PostgreSQL.

\noindent\begin{minipage}{\linewidth}
\lstinputlisting[
    style=evidence,firstline=8,lastline=12,
    caption={Captured dimension loading from Silver into PostgreSQL.},
    label={lst:batch-dimensions}
]{\evidenceroot/02-dimensions.txt}
\end{minipage}


\noindent\begin{minipage}{\linewidth}
\lstinputlisting[
    style=evidence,firstline=8,lastline=15,
    caption={Captured Spark Gold aggregation and database publication.},
    label={lst:batch-gold}
]{\evidenceroot/03-gold.txt}
\end{minipage}


The following result tables are database queries over the batch-produced Gold
tables, captured by \texttt{make answers}. They show sample answers to the first
two user questions. The accompanying times are individual query response times,
not batch processing durations.
% This run's Gold log predates the restored Spark DataFrame displays.
% These listings are explicitly database results, not DataFrame console output.

\paragraph{Question 1: Which states had the highest average PM2.5 in 2024?}
Arkansas has the highest annual average in this output (11.09).

\noindent\begin{minipage}{\linewidth}
\lstinputlisting[
    style=evidence,firstline=13,lastline=31,
    caption={Top ten state-year averages queried from the batch Gold results.},
    label={lst:batch-answer-states}
]{\evidenceroot/06-answers.txt}
\end{minipage}


\paragraph{Question 2: Which county-days had the highest average PM2.5?}
The highest county-day average shown is 217.32 for Fairbanks North Star,
Alaska, on June 30, 2024.

\noindent\begin{minipage}{\linewidth}
\lstinputlisting[
    style=evidence,firstline=34,lastline=52,
    caption={Top ten county-day averages queried from the batch Gold results.},
    label={lst:batch-answer-county-days}
]{\evidenceroot/06-answers.txt}
\end{minipage}


\subsection{Part 2: Stream}
% Show evidence of event-time-ordered Kafka replay and a running consumer result.
% Include sample producer/consumer outputs.

\subsection{Part 3: Serve}
% Show that batch and stream results are loaded into the chosen database.
% Supply each user question, its database answer, and measured response time.
% Expand the table or use separate query/output blocks as needed.

\begin{table}[htbp]
    \centering
    \caption{User questions, database answers, and response times.}
    \begin{tabularx}{\linewidth}{@{}XXX@{}}
        \toprule
        User question & Answer / output & Response time (ms) \\
        \midrule
        & & \\
        & & \\
        & & \\
        \bottomrule
    \end{tabularx}
\end{table}

\subsection{Part 4: Quality and Monitoring}
% Include at least four automated pipeline checks and failed-row counts.
% Show sample check outputs and rows in/out at every pipeline stage.

\begin{table}[htbp]
    \centering
    \caption{Automated quality-check results.}
    \begin{tabularx}{\linewidth}{@{}XXr@{}}
        \toprule
        Check & Failure criterion & Failed rows \\
        \midrule
        & & \\
        & & \\
        & & \\
        & & \\
        \bottomrule
    \end{tabularx}
\end{table}

\begin{table}[htbp]
    \centering
    \caption{Row counts at each pipeline stage.}
    \begin{tabularx}{\linewidth}{@{}Xrr@{}}
        \toprule
        Pipeline stage & Rows in & Rows out \\
        \midrule
        & & \\
        & & \\
        & & \\
        & & \\
        \bottomrule
    \end{tabularx}
\end{table}

\section{The Deep Part}
% 3 marks.

\subsection{What Changed}
% Identify the selected area, the baseline, and the change implemented.

\subsection{Measurements Before and After}
% Describe the measurement method and conditions; supply comparable results.
% Add metric names and units in the table.

\begin{table}[htbp]
    \centering
    \caption{Measurements before and after the change.}
    \begin{tabularx}{\linewidth}{@{}XXX@{}}
        \toprule
        Metric (unit) & Before & After \\
        \midrule
        & & \\
        & & \\
        & & \\
        \bottomrule
    \end{tabularx}
\end{table}

\subsection{What We Learned}
% Interpret the measurements and explain lessons, tradeoffs, and limitations.

\section{At 100 Times the Data}
% 1 mark.

\subsection{What Breaks First}
% Identify the first bottleneck and explain why it would fail at this scale.

\subsection{What We Would Change}
% Explain the changes needed to address the bottleneck.

\section{Contributions and AI Acknowledgement}
% 1 mark.

\subsection{Contribution Table}
% Add one row for each team member and describe who did what.

\begin{table}[htbp]
    \centering
    \caption{Team contributions.}
    \begin{tabularx}{\linewidth}{@{}lX@{}}
        \toprule
        Team member & Contributions \\
        \midrule
        & \\
        & \\
        & \\
        & \\
        \bottomrule
    \end{tabularx}
\end{table}

\subsection{AI Acknowledgement}
% Name each AI tool or outside-code source and state what it was used for,
% as required by the project outline. Include assistance with this template.

% Optional appendix: uncomment only if needed.
% Keep the main report within eight pages.
% \clearpage
% \appendix
% \section{Supporting Material}
% Add supplementary outputs, logs, queries, or measurement details.

\end{document}
